\section{Finite calibration, scoring and extreme-confidence benchmarks}
\label{sec:annex-p19-calibration-extreme-confidence}

\begin{annexitem}{annex:p19.calibration-decomposition}{thm:p19-calibration-decomposition}{Conditional squared error and recalibration}
\begin{proof}
All sums are finite. On a judgment class \(C=c\) of positive mass, write
\(m=m(c)\). Since \(Y\) is binary, \(\mathbb E[Y\mid C=c]=m\)
and \(\mathbb E[Y^2\mid C=c]=m\). Thus for any real \(a\),
\[
 \mathbb E[(a-Y)^2\mid C=c]
 =a^2-2am+m=(a-m)^2+m(1-m).
\]
Take \(a=g(c)\), multiply by \(w_c\), and sum over positive-mass
classes. The grouping identity of \cref{thm:p15-finite-total-expectation}
gives the asserted decomposition. Choosing \(g(c)=c\) gives the first
formula for \(L\). Choosing \(g(c)=m(c)\) eliminates the nonnegative
square term; equality in the minimization requires \(g(c)=m(c)\) on
every positive-mass class. This proves almost-sure uniqueness and reduction
of the score by \(\operatorname{Rel}\).

Total expectation gives \(\sum_c w_cm(c)=\mu\). Consequently
\[
 \mathbb E[m(C)(1-m(C))]
 =\mu-\mathbb E[m(C)^2]
 =\mu(1-\mu)-\mathbb E[(m(C)-\mu)^2].
\]
This proves the second decomposition. Since \(0\leq m(c)\leq1\),
the first expectation is nonnegative, giving the upper bound on resolution;
its lower bound follows from its squared definition.

For a value \(q\) of \(Q=m(C)\) with positive probability, group the
original classes with \(m(c)=q\):
\[
 \mathbb P(Y=1,Q=q)=\sum_{c:m(c)=q}w_cm(c)
 =q\sum_{c:m(c)=q}w_c=q\mathbb P(Q=q).
\]
Divide by the positive class mass. This proves calibration even if the map
\(m\) merges classes. Calibration of \(C\) itself implies
\(\mathbb E[C]=\mathbb E[m(C)]=\mu\). The explicit counterexample
to the converse is calculated in \cref{annex:p19.score-calculations}.
No independence among repeated observations is needed for these identities.
\end{proof}
\end{annexitem}

\begin{annexitem}{annex:p19.binned-calibration}{thm:p19-binned-calibration}{Within-bin variance and covariance terms}
\begin{proof}
Fix a bin of positive probability, with means \(a,b\), and expand
\(C-Y=(a-b)+(C-a)-(Y-b)\). The conditional means of \(C-a\) and
\(Y-b\) are zero, eliminating their cross terms with \(a-b\). Thus
\[
 \mathbb E[(C-Y)^2\mid B]
 =(a_B-b_B)^2+v_B+b_B(1-b_B)-2\gamma_B.
\]
The conditional variance of binary \(Y\) is \(b_B(1-b_B)\), by the
two-value calculation in the preceding proof. Average over bins and use
\(\mathbb E[b_B]=\mu\) to obtain
\[
 \mathbb E[b_B(1-b_B)]
 =\mu(1-\mu)-\mathbb E[(b_B-\mu)^2].
\]
This proves the identity. If \(C\) is constant in each bin, then
\(C-a_B=0\), hence \(v_B=\gamma_B=0\).

For failure of the converse, put both classes of the miscalibrated law in
\cref{def:p19-brier-calibration-examples} in one bin. Both bin means
equal .75, while the reports .6 and .9 have conditional outcome rates
.9 and .6. Explicitly, \(v_B=.0225\) and
\[
 \gamma_B=\tfrac12(.6-.75)(.9-.75)
          +\tfrac12(.9-.75)(.6-.75)=-.0225.
\]
The binned identity gives \(.1875+.0225+2(.0225)=.255\).
Omitting the two additional terms would give the wrong original score.
\end{proof}
\end{annexitem}

\begin{annexitem}{annex:p19.brier-propriety}{thm:p19-brier-propriety}{The original categorical score and binary normalization}
\begin{proof}
For a probability vector \(f\) and realized category \(k\), the loss is
\[
 \sum_j(f_j-\mathbf1_{\{j=k\}})^2=\sum_jf_j^2-2f_k+1.
\]
It is nonnegative as a sum of squares. Nonnegative coordinates summing to
one have \(\sum_jf_j^2\leq1\), so the loss is at most two. A unit
report on the realized category gives zero; a unit report on another gives
two. Another category exists for \(r\geq2\). The same range holds
for empirical averages.

Average over the outcome law \(q\):
\[
 \sum_jf_j^2-2\sum_jf_jq_j+1
 =\sum_j(f_j-q_j)^2+1-\sum_jq_j^2.
\]
The squared term vanishes precisely at the admissible report \(f=q\).
This proves strict propriety. Apply the same calculation to the uniform
law on the fixed empirical record, with \(q_j\) its category frequency,
to prove the best-constant-report result. Knowing this empirical optimum
after observing outcomes is not knowing future outcome frequencies.

For two categories the two squares are \((c-y)^2\) and
\(((1-c)-(1-y))^2=(y-c)^2\). Their sum is \(2(c-y)^2\), proving
the normalization relation without an independence premise.

To verify the body's distinction between belief and outcome law, suppose
a person assigns probability .9 to \(Y=1\) and minimizes expected
squared error under that belief. The result just proved uniquely selects
report .9. Under a law with \(\mathbb P(Y=1)=.6\), however, a constant
.9 report has \(m(.9)=.6\) and is not calibrated. Its actual loss is
\(.6(.1)^2+.4(.9)^2=.33\), whereas report .6 has loss
\(.6(.4)^2+.4(.6)^2=.24\). The latter is less extreme. Both reports
are constant, so an ordinary correlation coefficient dividing by forecast
standard deviation is undefined. This proves the counterexample to a
necessary increase in extremity or correlation whenever score improves.
\end{proof}
\end{annexitem}

\begin{annexitem}{annex:p19.score-calculations}{thm:p19-score-calculations}{The forecast table, constructed laws and odds bounds}
\begin{proof}
The squared errors of Brier's ten rain forecasts are
\(.49,.01,.04,.36,.04,0,0,0,0,.01\). They sum to .95, giving a
one-coordinate average .095 and original category-sum score .19.
Three of the ten occasions have rain, so a constant .3 report has
\[
 2\bigl(.3(.3-1)^2+.7(.3-0)^2\bigr)
 =2(.147+.063)=.42.
\]
A constant .2 gives \(2(.3\cdot.64+.7\cdot.04)=.44\). The optimal
constant result also follows from the category law \((.3,.7)\), with
\(1-(.3^2+.7^2)=.42\).

The two-class examples are actual finite probability laws: assign masses
\(m_j/2\) and \((1-m_j)/2\) to outcomes \(Y=1\) and \(Y=0\)
in class \(j=1,2\), and attach the stipulated report to that class.
All masses are nonnegative and sum to one. The constant example can
instead use two outcomes of masses .75 and .25. Its conditional mean
equals its report .75; both resolution and calibration error vanish,
and its loss is \(.75(1-.75)=.1875\).

For matching reports and rates \((.6,.9)\),
\(\mu=(.6+.9)/2=.75\) and
\[
 \operatorname{Res}=\tfrac12(.6-.75)^2+\tfrac12(.9-.75)^2=.0225.
\]
Calibration error is zero, so the decomposition gives .165. Interchanging
the rates retains the same mean judgment, mean outcome and resolution,
but calibration error is
\(\tfrac12(.6-.9)^2+\tfrac12(.9-.6)^2=.09\). The score is
therefore \(.09+.1875-.0225=.255\). Equal means coexist with
miscalibration on both positive-mass classes, proving the counterexample.
If both conditional means are instead .75 while the two reports remain
.6 and .9, then \(m(C)=\mu=.75\) is constant and resolution is zero.
Report variance and calibration error are both
\(\tfrac12(.6-.75)^2+\tfrac12(.9-.75)^2=.0225\), so the loss is
\(.0225+.1875=.21\). This proves that varying reported confidence
need not provide any resolution.

For \(p<1\), odds against the complementary event are \(p/(1-p)\).
Solving \(p/(1-p)=t\) gives \(p=t/(1+t)\). Conversely this probability
has those odds, because \(1-p=1/(1+t)>0\) for every finite \(t\geq0\).
Multiplying the ratio by 100 gives all listed percentage benchmarks.
In particular error percentages at odds 50, 100 and 1000 are exactly
\(100/51\), \(100/101\) and \(100/1001\), before rounding.

Under calibration the conditional error probability at report \(c\) is
\(1-c\). Thus
\[
 \mathbb P(Y=0,A_t)=\sum_{c\geq t/(1+t)}w_c(1-c)
 \leq\frac1{1+t}\sum_{c\geq t/(1+t)}w_c.
\]
Divide by the positive mass of \(A_t\). At \(C=1\), the error
contribution is zero. At \(C=0\), the probability of \(Y=1\) is
\(w_0m(0)=0\) if the class has positive mass, and is already zero
if it is null. This proves the endpoint identities. The finite observed
rate at a finite odds threshold is a separate random or recorded quantity,
so these population expectations impose no equality on every realization.
\end{proof}
\end{annexitem}

\begin{annexitem}{annex:p19.certainty-game}{thm:p19-certainty-game}{Linearity for the chip game and subjective expected money}
\begin{proof}
Let \(Z_i\) indicate a red chip on selected answer \(i\). Conditional
on the already scored answers, the unprotected net payment is
\(G=b\sum_{i=1}^kZ_i-w\). Each indicator has expectation
\(2/(100+2)=1/51\). By \cref{thm:p15-finite-expectation-laws},
\(\mathbb E[G]=kb/51-w\). Neither a factorization of answer
probabilities nor independent answers were used. Replacing the chip before
the next random draw maintains the stated marginal chip probability.

Before the answer outcomes are known, average under the participant's
conditional beliefs \(q_i\) for correctness and \(1/51\) for each
chip's red probability. Linearity gives
\[
 \mathbb E_{\mathrm{belief}}[G]=\frac{kb}{51}-\sum_{i=1}^k(1-q_i).
\]
If every \(q_i\geq50/51\), the error sum is at most \(k/51\).
At \(b=1\), the expected net is \(\sum_i(q_i-50/51)\); it is
positive precisely when some term is positive. Preferring larger expected
money is an additional decision criterion; the calculation asserts nothing
about every possible utility function or observed willingness to play.

For seventeen selected answers the expected chip count is \(17/51=1/3\).
Fourteen white and two red chips give red probability \(2/16=1/8\),
balancing an error probability of one eighth at equal stakes. These are
expectations, not identities for realized draw sequences.

Experiment~5 had no realized red chips among its players' selected answers,
so each unprotected net was minus that player's number of wrong answers.
The 46 reported errors give unprotected aggregate balance \(-46\) dollars
and error fraction \(46/387\approx.11886\). Subsequent protection of
the promised participation payment changes actual payments, not these
scored answers or the unprotected-balance calculation. Replacing realized
zero-red outcomes with expected chip counts produces a different quantity.
The paper does not give per-player selected counts and stakes sufficient
to reproduce each reported simulated balance independently.
\end{proof}
\end{annexitem}
